\documentclass[../main.tex]{subfiles}
\begin{document}

\section*{Supporting Information}
\label{sec:suppinfo}

\clearpage
\setcounter{figure}{0}
\renewcommand{\thefigure}{S\arabic{figure}}
\setcounter{table}{0}
\renewcommand{\thetable}{S\arabic{table}}
\setcounter{equation}{0}
\renewcommand{\theequation}{S\arabic{equation}}

% ─────────────────────────────────────────────────────────────────────────────
\subsection*{S1\quad Prompt Used for Single-Agent RoB Assessment}
\label{sec:si-prompt}

The following prompt was used without modification across all eight models and all
227~studies.
Each study was assessed in seven sequential inference calls---one per RoB domain---with
the full Markdown text of the study report included in every call as the user message,
followed by the domain identifier (e.g.\ \texttt{Domain: D1}).
The prompt text below was supplied as the system message.

\small
\begin{verbatim}
Goal: Systematically assess the risk of bias in a research article using
Cochrane RoB 1.0. Evaluate bias domain, assign a judgment (Low, High,
or Unclear Risk), and justify strictly from the article evidence.

---

Preliminary Considerations

Before assessing, identify:
1. Intervention Type: Pharmacological | Surgical/Procedural |
   Non-pharmacological | Device
2. Primary Outcome Type: Objective (mortality, lab values) |
   Subjective (pain, QoL)
3. Comparator: Placebo | Active comparator | Usual care | No intervention

Critical Context: Blinding expectations vary by intervention type.
For pharmacological interventions, lack of placebo is a serious flaw.
For surgical/non-pharmacological interventions where blinding is impossible,
focus on compensatory measures: blinding of outcome assessors and objective
endpoints.

---

Bias Domains

1. Random Sequence Generation (Selection Bias)
   Signalling Question: Was a truly random method used?
   Low Risk:    Random method adequately described (computer-generated,
                random tables, coin toss, dice, shuffled envelopes,
                minimization with random element)
   High Risk:   Non-random method (alternation, date of birth, hospital
                record number, clinician judgment)
   Unclear Risk: States "randomized" but no method details provided

2. Allocation Concealment (Selection Bias)
   Signalling Question: Was the allocation sequence concealed from those
   enrolling participants until assignment?
   Low Risk:    Adequate concealment (central/telephone randomization,
                sequentially numbered opaque sealed envelopes, numbered
                drug containers, independent pharmacy)
   High Risk:   Allocation not concealed or predictable (open list,
                non-opaque envelopes, alternation)
   Unclear Risk: Insufficient information to judge
   Note: Allocation concealment (before randomization) is different from
   blinding (after randomization).

3. Blinding of Participants and Personnel (Performance Bias)
   Signalling Questions: Were participants/personnel aware of assignments?
   If yes, could this affect outcomes?
   Pharmacological: Double-blinding expected; no placebo implies higher
   risk for subjective outcomes.
   Surgical/Non-pharmacological: Blinding may be impossible and is not
   automatically High Risk; check compensatory safeguards.
   Objective outcomes: Less susceptible to performance bias.
   Low Risk:    Adequate blinding, OR blinding impossible with compensatory
                measures, OR outcome is objective/hard endpoint
   High Risk:   No blinding + subjective outcome + no compensatory measures
   Unclear Risk: Insufficient blinding information

4. Blinding of Outcome Assessors (Detection Bias)
   Signalling Questions: Were assessors blinded? Is outcome objective or
   subjective?
   Low Risk:    Assessors blinded, OR objective outcome unlikely to be
                influenced by assignment knowledge
   High Risk:   Assessors not blinded + subjective outcome requiring
                interpretation
   Unclear Risk: Insufficient information
   Note: For patient-reported outcomes, participant acts as assessor.

5. Incomplete Outcome Data (Attrition Bias)
   Signalling Questions: Are data available for most participants? Are
   missing-data reasons balanced/unrelated to outcome?
   Low Risk:    Minimal missingness (<5%), OR balanced missingness with
                unrelated reasons, OR valid ITT/imputation
   High Risk:   High missingness (>10-20%), OR imbalance between groups,
                OR outcome-related missingness, OR per-protocol exclusions
   Unclear Risk: Missingness not adequately reported

6. Selective Reporting (Reporting Bias)
   Signalling Questions: Is protocol/registry available? Were pre-specified
   outcomes reported?
   Low Risk:    All pre-specified outcomes reported, OR protocol/registry
                aligns with publication
   High Risk:   Selective non-reporting (methods outcomes missing in
                results, changed primary outcome, significance-driven
                reporting)
   Unclear Risk: No protocol/registry and insufficient information

7. Other Bias
   Signalling Questions: Are there other bias sources (early stopping,
   baseline imbalance, conflicts, contamination)?
   Low Risk:    No relevant additional bias detected
   High Risk:   Important additional bias present
   Unclear Risk: Insufficient information

---

Key Principles
1. Evidence-based only: Base judgments on explicit information. Do not
   assume. If evidence is missing, choose "Unclear Risk".
2. Cite evidence: Reference specific text or note absence of information.
3. Context matters: Apply blinding expectations appropriate to intervention
   type.
4. Impossibility vs Negligence: Not blinding when impossible differs from
   not blinding when feasible.
5. Outcome sensitivity: Subjective outcomes require stricter blinding than
   objective ones.

---

Return ONLY valid JSON with:
{
  "judgement": "Low Risk" | "High Risk" | "Unclear Risk",
  "justification": "short justification"
}
\end{verbatim}
\normalsize

% ─────────────────────────────────────────────────────────────────────────────
\subsection*{S2\quad Overall Model Performance}
\label{sec:si-overall}

\begin{sidewaystable}
\centering
\caption{Overall performance on RoB~1.0 assessment ($n = 227$ studies;
1{,}589 cells). Models ranked by accuracy.
F1-LR, F1-HR, F1-UR: F1 for Low Risk, High Risk, and Unclear Risk computed
globally across all study--domain cells.
Cohen's~$\kappa$ and Gwet's~AC1 are both chance-corrected agreement
coefficients; they differ in how the expected chance agreement $p_e$ is
computed (see Methods).
Best value in each column shown in bold.
95\% confidence intervals (study-clustered bootstrap, 5{,}000 iterations,
resampling 227 studies with replacement) are shown in \textit{brackets} below
each point estimate; all analyses are exploratory (see Methods, Statistical
Inference).}
\label{tab:overall}
\begin{tabular}{lccccccc}
\toprule
Model & Accuracy (\%) & Macro-F1 & F1-LR & F1-HR & F1-UR & Cohen's $\kappa$ & Gwet's AC1 \\
\midrule
Claude~Haiku~4.5   & \textbf{62.5}   & 0.541           & \textbf{0.743}  & 0.469           & 0.413           & 0.325           & \textbf{0.485}  \\
                   & \footnotesize[59.8, 65.2] & \footnotesize[0.511, 0.569] & \footnotesize[0.718, 0.767] & \footnotesize[0.412, 0.521] & \footnotesize[0.361, 0.463] & \footnotesize[0.286, 0.365] & \footnotesize[0.443, 0.526] \\
\addlinespace[3pt]
GLM-5              & 61.1            & 0.557           & 0.709           & 0.441           & \textbf{0.520}  & 0.341           & 0.449           \\
                   & \footnotesize[58.7, 63.5] & \footnotesize[0.529, 0.582] & \footnotesize[0.682, 0.735] & \footnotesize[0.385, 0.494] & \footnotesize[0.479, 0.561] & \footnotesize[0.301, 0.380] & \footnotesize[0.411, 0.485] \\
\addlinespace[3pt]
Qwen~3.7-Plus      & 60.9            & 0.525           & 0.718           & 0.356           & 0.503           & 0.306           & 0.457           \\
                   & \footnotesize[58.3, 63.4] & \footnotesize[0.496, 0.552] & \footnotesize[0.691, 0.743] & \footnotesize[0.293, 0.413] & \footnotesize[0.459, 0.543] & \footnotesize[0.266, 0.345] & \footnotesize[0.417, 0.495] \\
\addlinespace[3pt]
Gemini~2.5~Flash   & 60.7            & \textbf{0.558}  & 0.718           & \textbf{0.474}  & 0.484           & \textbf{0.349}  & 0.437           \\
                   & \footnotesize[58.2, 63.1] & \footnotesize[0.533, 0.582] & \footnotesize[0.690, 0.744] & \footnotesize[0.424, 0.516] & \footnotesize[0.443, 0.526] & \footnotesize[0.312, 0.384] & \footnotesize[0.399, 0.475] \\
\addlinespace[3pt]
DeepSeek~v4-Flash  & 60.3            & \textbf{0.558}  & 0.704           & \textbf{0.480}  & 0.489           & 0.337           & 0.433           \\
                   & \footnotesize[57.7, 62.9] & \footnotesize[0.531, 0.585] & \footnotesize[0.676, 0.730] & \footnotesize[0.429, 0.530] & \footnotesize[0.447, 0.530] & \footnotesize[0.296, 0.377] & \footnotesize[0.394, 0.472] \\
\addlinespace[3pt]
Grok~4.3           & 60.1            & 0.549           & 0.696           & 0.432           & 0.518           & 0.326           & 0.433           \\
                   & \footnotesize[57.8, 62.5] & \footnotesize[0.523, 0.574] & \footnotesize[0.670, 0.722] & \footnotesize[0.376, 0.486] & \footnotesize[0.478, 0.556] & \footnotesize[0.288, 0.362] & \footnotesize[0.397, 0.471] \\
\addlinespace[3pt]
GPT-5.4            & 57.7            & 0.546           & 0.668           & 0.464           & 0.506           & 0.322           & 0.387           \\
                   & \footnotesize[55.1, 60.3] & \footnotesize[0.519, 0.572] & \footnotesize[0.636, 0.698] & \footnotesize[0.410, 0.511] & \footnotesize[0.468, 0.543] & \footnotesize[0.282, 0.361] & \footnotesize[0.349, 0.426] \\
\addlinespace[3pt]
Gemini~2.5~Pro     & 57.6            & 0.543           & 0.674           & 0.455           & 0.502           & 0.331           & 0.383           \\
                   & \footnotesize[55.1, 60.0] & \footnotesize[0.517, 0.569] & \footnotesize[0.643, 0.702] & \footnotesize[0.407, 0.501] & \footnotesize[0.458, 0.542] & \footnotesize[0.293, 0.367] & \footnotesize[0.344, 0.420] \\
\bottomrule
\end{tabular}
\end{sidewaystable}

% ─────────────────────────────────────────────────────────────────────────────
\subsection*{S3\quad Per-Domain Performance by Model}
\label{sec:si-domain}

Tables~\ref{tab:domain-full} and~\ref{tab:domain-f1} report accuracy (\%) and
Macro-F1 for each model in each of the seven RoB domains.

\begin{sidewaystable}
\centering
\caption{Per-domain accuracy (\%) for all eight models
($n = 227$ cells per domain). Best value per domain in bold.}
\label{tab:domain-full}
\begin{tabular}{lC{1.5cm}C{1.5cm}C{1.5cm}C{1.5cm}C{1.5cm}C{1.5cm}C{1.5cm}}
\toprule
Model & D1 & D2 & D3 & D4 & D5 & D6 & D7 \\
\midrule
Claude~Haiku~4.5   & 73.6 & 74.0 & 67.4          & 51.5 & \textbf{64.3} & 57.7 & 48.9 \\
GLM-5              & 76.7 & 74.0 & 55.9          & 53.7 & 62.1          & 55.9 & 49.3 \\
Qwen~3.7-Plus      & \textbf{77.5} & 74.9 & 49.8 & 52.0 & 59.9 & 56.4 & \textbf{55.5} \\
Gemini~2.5~Flash   & 72.2 & 74.0 & \textbf{72.2} & 51.5 & 60.4 & \textbf{62.6} & 31.7 \\
DeepSeek~v4-Flash  & 74.0 & 70.0 & 67.0          & \textbf{57.7} & 54.6 & 55.9 & 42.7 \\
Grok~4.3           & 75.8 & 70.0 & 59.5          & 54.6 & 58.6 & 60.8 & 41.4 \\
GPT-5.4            & 75.8 & 70.0 & 64.3          & 52.0 & 51.5 & 55.9 & 34.4 \\
Gemini~2.5~Pro     & 76.7 & \textbf{74.4} & 62.6  & 55.1 & 55.1 & 55.1 & 24.2 \\
\bottomrule
\end{tabular}
\end{sidewaystable}

\begin{sidewaystable}
\centering
\caption{Per-domain Macro-F1 for all eight models
($n = 227$ cells per domain). Best value per domain in bold.}
\label{tab:domain-f1}
\begin{tabular}{lC{1.5cm}C{1.5cm}C{1.5cm}C{1.5cm}C{1.5cm}C{1.5cm}C{1.5cm}}
\toprule
Model & D1 & D2 & D3 & D4 & D5 & D6 & D7 \\
\midrule
Claude~Haiku~4.5   & 0.572 & 0.618 & 0.527          & 0.425 & 0.394 & 0.387 & 0.269 \\
GLM-5              & 0.632 & 0.657 & 0.462          & 0.479 & \textbf{0.514} & 0.443 & 0.345 \\
Qwen~3.7-Plus      & \textbf{0.638} & 0.672 & 0.428 & 0.447 & 0.439 & 0.388 & 0.288 \\
Gemini~2.5~Flash   & 0.549 & 0.663 & \textbf{0.626} & 0.411 & 0.412 & 0.438 & 0.280 \\
DeepSeek~v4-Flash  & 0.602 & 0.610 & 0.583          & \textbf{0.521} & 0.419 & 0.424 & \textbf{0.347} \\
Grok~4.3           & 0.618 & 0.668 & 0.524          & 0.492 & 0.423 & \textbf{0.449} & 0.315 \\
GPT-5.4            & 0.625 & 0.635 & 0.515          & 0.469 & 0.419 & 0.397 & 0.317 \\
Gemini~2.5~Pro     & 0.621 & \textbf{0.671} & 0.481  & 0.478 & 0.413 & 0.475 & 0.195 \\
\bottomrule
\end{tabular}
\end{sidewaystable}

% ─────────────────────────────────────────────────────────────────────────────
\subsection*{S4\quad Stratification by Intervention Type}
\label{sec:si-itype}

Figure~\ref{fig:si-cdomain-itype} shows the distribution of intervention types
within each clinical domain in the 227-study corpus.
Because each study can be tagged with multiple clinical domains, the counts
across domains sum to more than 227; a study tagged with two domains contributes
its intervention type to both.
The distribution is markedly non-uniform: Pregnancy \& childbirth and
Complementary \& alternative medicine are almost exclusively pharmacological,
Effective practice \& health systems and Mental health are predominantly
non-pharmacological, and Orthopaedics \& trauma contains no pharmacological
studies.
Table~\ref{tab:itype-full} reports Macro-F1 for each model within each intervention
type. Figure~\ref{fig:si-itype} shows the same data as a heatmap.

\begin{figure}[H]
  \centering
  \includegraphics[width=\linewidth]{all_reviews_clinical_domain_intervention_table}
  \caption{Distribution of intervention types across clinical domains
  ($n = 227$ studies; studies with multiple clinical domain tags are counted
  in each relevant domain).
  Cell colour encodes the fraction of that domain's studies belonging to each
  intervention type (pastel red = low proportion, pastel blue = high proportion);
  grey cells indicate zero studies.
  Domains are sorted by total study count (descending).}
  \label{fig:si-cdomain-itype}
\end{figure}

\begin{table}[ht]
\centering
\caption{Macro-F1 stratified by intervention type. Numbers in parentheses are
study counts ($n_{\text{studies}}$). Best value per row in bold.
Subgroup estimates carry substantial uncertainty given the small group sizes;
these analyses are descriptive only.}
\label{tab:itype-full}
\begin{tabular}{lcccc}
\toprule
Model
  & \makecell{Drug /\\medication\\($n=104$)}
  & \makecell{Non-pharma-\\cological\\($n=75$)}
  & \makecell{Procedure\\($n=30$)}
  & \makecell{Medical\\device\\($n=18$)} \\
\midrule
Claude~Haiku~4.5   & 0.548 & 0.530 & 0.494 & \textbf{0.627} \\
GLM-5              & \textbf{0.559} & \textbf{0.590} & 0.514 & 0.471 \\
Qwen~3.7-Plus      & \textbf{0.511} & 0.567 & 0.474 & 0.445 \\
Gemini~2.5~Flash   & 0.540 & 0.574 & 0.529 & \textbf{0.596} \\
DeepSeek~v4-Flash  & 0.548 & \textbf{0.582} & \textbf{0.520} & 0.543 \\
Grok~4.3           & \textbf{0.552} & 0.561 & 0.519 & 0.531 \\
GPT-5.4            & 0.515 & \textbf{0.595} & \textbf{0.530} & 0.502 \\
Gemini~2.5~Pro     & 0.533 & 0.575 & 0.498 & 0.534 \\
\bottomrule
\end{tabular}
\end{table}

\begin{figure}[ht]
  \centering
  \includegraphics[width=\linewidth]{all_reviews_intervention_type_f1_heatmap}
  \caption{Macro-F1 heatmap by intervention type and model.
  Colour scale: 0--1.}
  \label{fig:si-itype}
\end{figure}

% ─────────────────────────────────────────────────────────────────────────────
\subsection*{S5\quad Stratification by Clinical Domain}
\label{sec:si-clinical}

Table~\ref{tab:clinical-full} reports mean Macro-F1 (averaged across all eight models)
for all 24 Cochrane Review Group categories, ordered from highest to lowest mean~F1.
Domains with fewer than five studies are included for completeness but should be
interpreted with caution.

\begin{table}[ht]
\centering
\caption{Mean Macro-F1 across all eight models, stratified by clinical domain
(24 Cochrane Review Groups, multi-label; study counts sum to more than 227).
Ordered by descending mean F1.
Estimates for smaller groups carry substantial uncertainty; these analyses are
descriptive only.}
\label{tab:clinical-full}
\begin{tabular}{lcc}
\toprule
Clinical domain & $n_{\text{studies}}$ & Mean Macro-F1 \\
\midrule
Orthopaedics \& trauma                 &  7 & 0.680 \\
Endocrine \& metabolic                 & 10 & 0.638 \\
Rheumatology                           &  5 & 0.637 \\
Wounds                                 &  5 & 0.628 \\
Dentistry \& oral health               & 15 & 0.597 \\
Lungs \& airways                       & 14 & 0.589 \\
Pregnancy \& childbirth                &  9 & 0.583 \\
Neonatal                               & 17 & 0.578 \\
Complementary \& alternative medicine  & 19 & 0.555 \\
Child health                           & 72 & 0.550 \\
Gastroenterology \& hepatology         & 21 & 0.538 \\
Heart \& circulation                   & 45 & 0.536 \\
Neurology                              & 11 & 0.519 \\
Cancer                                 & 21 & 0.494 \\
Effective practice \& health systems   &  4 & 0.487 \\
Kidney disease                         & 31 & 0.484 \\
Mental health                          & 11 & 0.481 \\
Gynaecology                            & 11 & 0.480 \\
Insurance medicine                     & 21 & 0.477 \\
Blood disorders                        &  6 & 0.476 \\
Genetic disorders                      &  9 & 0.473 \\
Skin disorders                         &  3 & 0.448 \\
Reproductive \& sexual health          &  8 & 0.437 \\
Infectious disease                     &  8 & 0.425 \\
\bottomrule
\end{tabular}
\end{table}

% ─────────────────────────────────────────────────────────────────────────────
\subsection*{S6\quad Pairwise Model Complementarity}
\label{sec:si-pairwise}

Table~\ref{tab:pairwise} reports pairwise complementarity across the eight models.
Each cell~$(i,j)$ gives the percentage of the 1{,}589 study--domain cells in which
model~$i$ produced an incorrect judgement and model~$j$ produced the correct one.
High values indicate that model~$j$ recovers errors that model~$i$ misses, i.e.\
that $j$ adds non-redundant information over $i$.
The matrix is asymmetric: cell~$(i,j) \neq$ cell~$(j,i)$ because the two models may
have different overall accuracy.

\begin{sidewaystable}
\centering
\caption{Pairwise complementarity matrix. Cell~$(i,j)$ = \% of 1{,}589 cells
where row model $i$ is incorrect and column model $j$ is correct.
Row model abbreviations: Haiku~=~Claude~Haiku~4.5; GLM-5~=~GLM-5;
Qwen~=~Qwen~3.7-Plus; Flash~=~Gemini~2.5~Flash; DSv4~=~DeepSeek~v4-Flash;
Grok~=~Grok~4.3; GPT~=~GPT-5.4; Pro~=~Gemini~2.5~Pro.}
\label{tab:pairwise}
\begin{tabular}{lcccccccc}
\toprule
 & \textbf{Haiku} & \textbf{GLM-5} & \textbf{Qwen} & \textbf{Flash} & \textbf{DSv4} & \textbf{Grok} & \textbf{GPT} & \textbf{Pro} \\
\midrule
Haiku &  ---  & 12.5 & 10.6 &  9.9 & 12.3 & 12.5 & 12.6 & 12.4 \\
GLM-5 & 13.9  &  ---  &  7.4 & 12.3 &  9.6 &  8.9 &  9.5 &  7.6 \\
Qwen  & 12.3  &  7.6  &  ---  & 11.7 & 10.3 &  7.6 &  9.6 &  9.4 \\
Flash & 11.7  & 12.8  & 11.9 &  ---  & 11.3 & 10.2 &  9.4 &  9.1 \\
DSv4  & 14.5  & 10.4  & 10.8 & 11.7 &  ---  &  9.6 &  9.5 &  9.9 \\
Grok  & 14.9  &  9.9  &  8.4 & 10.8 &  9.8 &  ---  &  9.0 &  8.4 \\
GPT   & 17.4  & 12.9  & 12.7 & 12.4 & 12.1 & 11.4 &  ---  & 10.1 \\
Pro   & 17.3  & 11.1  & 12.6 & 12.2 & 12.6 & 11.0 & 10.3 &  ---  \\
\bottomrule
\end{tabular}
\end{sidewaystable}

\subsection*{S7\quad Confidence Interval Method Comparison}
\label{sec:si-ci}

All confidence intervals reported in the main text were computed using
study-clustered bootstrap (5{,}000 iterations, resampling the 227 studies
with replacement).
This section documents the comparison between bootstrap and two Wilson score
interval variants for accuracy, and explains why the bootstrap approach was
adopted for all metrics.

\subsubsection*{Wilson score interval}

For a proportion $p$ observed over $n$ independent Bernoulli trials,
the Wilson 95\% CI is
\begin{equation}
  \left(
    \frac{n p + \tfrac{z^2}{2}}{n + z^2}
    \;\pm\;
    \frac{z\,\sqrt{n p(1-p) + \tfrac{z^2}{4}}}{n + z^2}
  \right),
  \label{eq:wilson}
\end{equation}
with $z = 1.96$.
Applied to accuracy ($p \approx 0.60$), two natural choices for $n$ exist:
$n = 1{,}589$ (all annotated cells, assuming independence) or
$n = 227$ (studies as the primary sampling unit).

\subsubsection*{Intraclass correlation and effective sample size}

The 1{,}589 cells are not independent: seven domain judgements are drawn from
the same study.
Pearson's one-way intraclass correlation coefficient (ICC) for the
correct/wrong binary indicator across the seven domains within each study
was estimated at $\hat\rho = 0.044$ (computed for Claude~Haiku~4.5 as a
representative model).
The low ICC indicates that whether a model correctly judges one domain bears
little information about its judgement on another domain of the same study,
which is consistent with the domain-specific nature of each RoB criterion.

Using the Kish~(1965) design-effect formula, the effective sample size is
\begin{equation}
  n_{\text{eff}} = \frac{n_{\text{cells}}}{1 + (k-1)\hat\rho}
                = \frac{1{,}589}{1 + 6 \times 0.044} \approx 1{,}257,
  \label{eq:neff}
\end{equation}
where $k = 7$ domains per study.

\subsubsection*{Numerical comparison}

Table~\ref{tab:ci-comparison} reports accuracy CIs for five models under three
methods.

\begin{table}[ht]
\centering
\caption{Comparison of 95\% confidence interval methods for overall accuracy.
Width in percentage points (pp).
Bootstrap: study-clustered, 5{,}000 iterations.
Wilson~$n_{\text{cells}}$: Equation~\eqref{eq:wilson} with $n = 1{,}589$.
Wilson~$n_{\text{eff}}$: Equation~\eqref{eq:wilson} with $n = 1{,}257$ (Kish formula).
Wilson~$n_{\text{studies}}$: Equation~\eqref{eq:wilson} with $n = 227$.}
\label{tab:ci-comparison}
\begin{tabular}{lcccc}
\toprule
Model & Accuracy & Wilson $n_{\text{cells}}$ (width) & Wilson $n_{\text{eff}}$ (width) & Bootstrap (width) \\
\midrule
Claude~Haiku~4.5  & 62.5\% & [60.1, 64.8] (4.8) & [59.8, 65.1] (5.3) & [59.8, 65.2] (5.4) \\
Gemini~2.5~Flash  & 60.7\% & [58.2, 63.0] (4.8) & [57.9, 63.4] (5.4) & [58.2, 63.1] (4.9) \\
DeepSeek~v4-Flash & 60.3\% & [57.9, 62.7] (4.8) & [57.6, 63.0] (5.4) & [57.7, 62.9] (5.2) \\
GPT-5.4           & 57.7\% & [55.3, 60.1] (4.9) & [54.9, 60.4] (5.4) & [55.1, 60.3] (5.2) \\
Gemini~2.5~Pro    & 57.6\% & [55.1, 60.0] (4.9) & [54.8, 60.3] (5.4) & [55.1, 60.0] (4.9) \\
\midrule
\multicolumn{2}{l}{Wilson $n_{\text{studies}} = 227$} & \multicolumn{3}{l}{widths $\approx$ 12.5--12.8 pp (not shown; see text)} \\
\bottomrule
\end{tabular}
\end{table}

Wilson with $n_{\text{cells}}$ underestimates uncertainty by $\sim$10\% relative
to bootstrap (widths 4.8~pp vs.\ 5.1~pp).
Wilson with $n_{\text{eff}}$ virtually matches the bootstrap (widths 5.3~pp
vs.\ 5.1~pp); the two methods agree to within 0.1~pp for every model.
Wilson with $n_{\text{studies}}$ overestimates uncertainty by $\sim$2.5$\times$
because it discards the domain-level information and treats each study as a
single binary observation.

\subsubsection*{Why bootstrap was adopted for all metrics}

The numerical equivalence between Wilson~$n_{\text{eff}}$ and bootstrap for
accuracy confirms that the low ICC ($\hat\rho = 0.044$) makes the independence
assumption nearly tenable.
Nevertheless, bootstrap was used uniformly for three reasons.
First, Macro-F1 and Gwet's~AC1 are composite statistics with no closed-form
variance expression; bootstrap applies without modification.
Second, the paired-bootstrap procedure used for inference-mode $\Delta$ values
(Table~\ref{tab:inference}) jointly resamples both conditions, correctly
propagating correlation between them; Wilson intervals cannot be directly
adapted for paired differences of non-proportion statistics.
Third, a single method for all reported CIs ensures internal consistency and
simplifies interpretation.


% ─────────────────────────────────────────────────────────────────────────────
\clearpage
\subsection*{S8\quad Inference Cost}
\label{sec:si-cost}

Table~\ref{tab:cost} reports total billed inference costs for the eight models
across the 227-study benchmark, retrieved from the OpenRouter account history.
All costs correspond to per-domain inference mode (seven sequential calls per study).
The total cost of running the full eight-model benchmark was approximately
\$244~USD, providing an order-of-magnitude reference for the practical
feasibility of this type of evaluation.

These figures should not be used to draw cross-model cost-efficiency conclusions.
The pipeline permitted up to two retries per domain call in the event of API
errors, timeouts, or invalid responses, and some studies required full
re-queuing; the number of retry attempts was not recorded per model.
Because billed costs reflect total API consumption including failed calls,
models with higher retry rates appear more expensive for reasons unrelated to
output quality or token generation.
Cross-model cost comparisons are therefore confounded by differential reliability
of provider APIs during the evaluation window and are not reported.

For Gemini~2.5~Flash and DeepSeek~v4-Flash, which were also evaluated in
single-call mode (Section~\ref{sec:inference}), the billed single-call costs
were \$2.05 and \$0.54 respectively, compared to \$9.42 and \$3.50 for
per-domain mode---a 4.6$\times$ and 6.5$\times$ reduction attributable to
the decrease from seven calls to one per study.
This cost differential is informative because both modes were run under
identical conditions for the same two models, making the within-model
comparison independent of the retry confound.

\begin{table}[ht]
\centering
\caption{Total billed inference costs for the eight models (per-domain mode,
227~studies), retrieved from OpenRouter account history, 7~July~2026.
Costs include retries due to API errors or invalid responses; the number of
retries was not recorded per model and cross-model comparisons should not be
drawn from these figures.
Total across all eight models: \$244.32.}
\label{tab:cost}
\begin{tabular}{lr}
\toprule
Model & Total billed cost (USD) \\
\midrule
DeepSeek~v4-Flash &  3.50 \\
Gemini~2.5~Flash  &  9.42 \\
Qwen~3.7-Plus     & 11.40 \\
GLM-5             & 30.70 \\
Claude~Haiku~4.5  & 32.40 \\
Grok~4.3          & 33.50 \\
Gemini~2.5~Pro    & 54.10 \\
GPT-5.4           & 69.30 \\
\midrule
Total             & 244.32 \\
\bottomrule
\end{tabular}
\end{table}

\end{document}
